\documentclass[10pt,a4paper]{amsart}
\usepackage[margin=19mm]{geometry}
\usepackage{amsmath,amssymb,mathtools}
\usepackage[scr=rsfs,cal=euler]{mathalfa}
\usepackage{xfrac}
\usepackage[unicode,hidelinks,bookmarks=false]{hyperref}
\newcommand{\ie}{i.e.\ }
\newcommand{\cf}{cf.\ }
\newcommand{\resp}{respectively\ }
\newcommand{\Subsection}{\subsection}
\providecommand{\nobreakdash}{\nobreak\hbox{-}}
\setlength{\emergencystretch}{2em}
\multlinegap0pt
\usepackage{xcolor}
\usepackage{amssymb}
\usepackage{csquotes}
\usepackage[shortlabels]{enumitem}
\usepackage{float}
\usepackage{mathtools}
\newtheorem{theorem}{Theorem}
\title{Spectral clustering via resistance Laplacian matrices}
\author{Priyavrat Deshpande, Gargi Lather}
\date{Source: arXiv:2609.04768v1 (CC BY 4.0)}
\begin{document}
\maketitle

\noindent\emph{Source and licence.} Spectral clustering via resistance Laplacian matrices. Version arXiv:2609.04768v1 (CC BY 4.0), distributed by arXiv under the Creative Commons Attribution 4.0 International licence, http://creativecommons.org/licenses/by/4.0/, as recorded in the arXiv licence metadata for this exact version and shown on the abstract page https://arxiv.org/abs/2609.04768v1. Full licence notice: \texttt{licenses/arxiv-2609.04768-CC-BY-4.0.txt}.\par\smallskip

\noindent\textit{Editorial scope.} Counted unit: the article's theorem thm:spectrumT (source0.tex lines 507--545). The article's complete original proof of it is delivered with it (source0.tex lines 547--699, one \texttt{proof} environment of 153 lines). No mathematics is written, repaired or re-derived: the statement and the proof are the article's own lines, in the article's own notation, and no retained line is substituted or re-flowed.  No further source-local context is needed: every cross-reference of every retained line resolves inside the two delivered spans.  Nothing was omitted from any retained span, so the package carries no omission record.  No retained line contains a citation, so the wrapper carries no bibliography and no bibliography entry.  No retained line contains a figure environment, an image-inclusion command or drawing code, so no image is delivered and the package declares no asset dependency.  Editorial formatting only, in addition: an AMS \texttt{amsart} preamble that restates the article's own declarations and macros used by the retained lines, namely the environments \texttt{theorem} on separate counters, together with the standard packages those lines need.  The retained environments print in this document as Theorem 1 in order of appearance, whereas the article prints the same items with the numbering of its own full document.  The complete native source of the article as distributed in the arXiv e-print for this exact version is bundled unchanged as \texttt{sources/209390/source0.tex}.
\begin{theorem}\label{thm:spectrumT}
Let $\alpha_1<\cdots<\alpha_k$ be the distinct resistance deviations,
where $\alpha_i$ occurs with multiplicity $m_i$, and let
\[
I_i=\{j\in\{1,\ldots,n\}:\delta_j=\alpha_i\}.
\]
Then $T\mathbf{1}=\mathbf{0}$. Moreover, for each $i$, the set
\[
E_i=
\left\{
x\in\mathbb{R}^n:
\operatorname{supp}(x)\subseteq I_i,\ 
\mathbf{1}^{\top}x=0
\right\}
\]
is precisely the set of solutions in $\mathbf{1}^{\perp}$ of
\[
Tx=\alpha_i x,
\]
and
\[
\dim E_i=m_i-1.
\]

The remaining $k-1$ eigenvalues of
$\left.T\right|_{\mathbf{1}^{\perp}}$ are simple and are precisely the
roots of
\[
f(\lambda)
=
\sum_{j=1}^{n}\frac{1}{\delta_j-\lambda}
=
\sum_{i=1}^{k}\frac{m_i}{\alpha_i-\lambda}.
\]
For each $i\in\{1,\ldots,k-1\}$, exactly one such root lies in
$(\alpha_i,\alpha_{i+1})$. Together with the eigenvalue $0$
corresponding to $\mathbf{1}$, these eigenvalues constitute the entire
spectrum of $T$.
\end{theorem}

\begin{proof}
By the identity $T=P\Delta P$, we have
$T\mathbf{1}=\mathbf{0}$. Since $T$ is symmetric,
$\mathbf{1}^{\perp}$ is invariant under $T$, and
\[
\mathbb{R}^n
=
\operatorname{span}\{\mathbf{1}\}
\oplus
\mathbf{1}^{\perp}
\]
is an orthogonal decomposition into $T$-invariant subspaces.

Let $x\in\mathbf{1}^{\perp}$ satisfy
\[
Tx=\lambda x.
\]
We obtain
\[
Tx
=
P\Delta x
=
\Delta x-\frac{1}{n}\mathbf{1}\delta^{\top}x.
\]
Set
\[
c=\frac{1}{n}\delta^{\top}x.
\]
The eigenvalue equation is then equivalent to
\[
(\delta_j-\lambda)x_j=c,
\qquad
j=1,\ldots,n.
\]

Suppose first that $c=0$. Then
\[
(\delta_j-\lambda)x_j=0
\]
for every $j$. Hence, if $x\neq 0$, then $\lambda=\alpha_i$ for some
$i$, and the support of $x$ is contained in $I_i$. Since
$x\in\mathbf{1}^{\perp}$, we have
$\mathbf{1}^{\top}x=0$, and therefore $x\in E_i$.

Conversely, if $x\in E_i$, then
\[
\delta^{\top}x
=
\alpha_i\mathbf{1}^{\top}x
=
0,
\]
and hence
\[
Tx
=
P\Delta x
=
\Delta x
=
\alpha_i x.
\]
Thus, $E_i$ is precisely the set of solutions in
$\mathbf{1}^{\perp}$ of $Tx=\alpha_i x$. Since the coordinates of a
vector in $E_i$ are supported on $I_i$ and have sum zero,
\[
\dim E_i=m_i-1.
\]

Suppose now that $c\neq 0$. Then
$\lambda\notin\{\alpha_1,\ldots,\alpha_k\}$, and
\[
x_j=\frac{c}{\delta_j-\lambda}.
\]
Since an eigenvector may be rescaled, we may assume that $c=1$.
The condition $x\in\mathbf{1}^{\perp}$ then becomes
\[
\sum_{j=1}^{n}\frac{1}{\delta_j-\lambda}=0.
\]

Conversely, suppose that
$\lambda\notin\{\alpha_1,\ldots,\alpha_k\}$ satisfies this equation,
and define
\[
x_j=\frac{1}{\delta_j-\lambda},
\qquad
j=1,\ldots,n.
\]
Then $\mathbf{1}^{\top}x=0$. Moreover,
\[
\frac{1}{n}\delta^{\top}x
=
\frac{1}{n}
\sum_{j=1}^{n}\frac{\delta_j}{\delta_j-\lambda}
=
\frac{1}{n}
\sum_{j=1}^{n}
\left(
1+\frac{\lambda}{\delta_j-\lambda}
\right)
=
1.
\]
Therefore,
\[
Tx
=
\Delta x-\mathbf{1}
=
\lambda x.
\]
Thus, the eigenvalues arising from the case $c\neq 0$ are precisely the
roots of $f$.

On each interval $(\alpha_i,\alpha_{i+1})$, the function $f$ is
continuously differentiable and satisfies
\[
f'(\lambda)
=
\sum_{j=1}^{n}\frac{1}{(\delta_j-\lambda)^2}
>0.
\]
Furthermore,
\[
\lim_{\lambda\to\alpha_i^+}f(\lambda)=-\infty,
\qquad
\lim_{\lambda\to\alpha_{i+1}^-}f(\lambda)=+\infty.
\]
Hence, $f$ has exactly one root in each interval
$(\alpha_i,\alpha_{i+1})$. Each root is simple because
$f'(\lambda)>0$.

For any such root $\lambda$, the relation
\[
x_j=\frac{c}{\delta_j-\lambda}
\]
shows that the corresponding eigenspace in $\mathbf{1}^{\perp}$ is
one-dimensional. Since $T$ is symmetric, $\lambda$ is a simple
eigenvalue of $\left.T\right|_{\mathbf{1}^{\perp}}$.

There are no roots in $(-\infty,\alpha_1)$, since $f(\lambda)>0$ on
that interval, and there are no roots in $(\alpha_k,\infty)$, since
$f(\lambda)<0$ there. Thus, $f$ has exactly $k-1$ roots.

Finally,
\[
\sum_{i=1}^{k}(m_i-1)+(k-1)=n-1.
\]
Therefore, the eigenvectors described above account for all of
$\mathbf{1}^{\perp}$. Together with the eigenvector $\mathbf{1}$,
they account for the entire spectrum of $T$.
\end{proof}

\end{document}
